\newpage
\section{Conclusion and Future Work}\label{sec:conclusion}

This thesis applies the complete AFML pipeline (CUSUM event filtering, triple-barrier labelling, fractional differentiation, sample weighting, CPCV with purging and embargo, the Deflated Sharpe Ratio, the Probability of Backtest Overfitting, and the DeLong pairwise AUC test) to BTC direction prediction on a 73-feature universe spanning four families: technical, mathematical, external (macroeconomic, crypto-macro, on-chain), and autoregressive lag. It benchmarks KAN against five comparison models (AR Logistic, Logistic Regression, Random Forest, XGBoost, LSTM) under identical CPCV splits ($N = 8$, $k = 2$, 28 splits, 7 backtest paths, 5 seeds per model), producing 840 prediction entries. It then extracts the first closed-form classification formula from a CPCV-trained KAN under the full AFML evaluation. The symbolic formula is the novel deliverable. The six-model benchmark is the methodological scaffolding that validates it.

Under leakage-free evaluation with multiple-testing correction, the literature's 80\% to 90\% accuracy claims do not survive. The top trained-model DSR (KAN at 0.2470) sits far below the 0.95 significance threshold, and $\text{PBO} = 0.657$ places the cross-section in the adversarial regime: 23 of 35 IS/OOS partitions see the in-sample best model underperform the out-of-sample median. The leave-one-out PBO identifies Random Forest as the anchor model whose removal leaves PBO unchanged. Six of seven bootstrap 95\% confidence intervals on median path Sharpe cross zero, and only buy-and-hold holds a strict-positive interval of $]1.8903, 2.4685[$. Buy-and-hold's median Sharpe (2.2722) exceeds every trained model on the same labelling window, at the cost of a $-99.98\%$ maximum drawdown. The honest answer to Q2 is that no architecture demonstrates a statistically significant predictive edge after correcting for selection bias. KAN posts the strongest single-model evidence (rank-1 trained median Sharpe and rank-1 trained DSR), but its bootstrap CI on the median Sharpe crosses zero. The DeLong test reaches 2 of 15 significant pairs at $\alpha = 0.05$ (AR Logistic versus XGBoost and Random Forest versus XGBoost, both with XGBoost as the lower-AUC half).

What survives is the symbolic-extraction contribution. The pipeline produces a humanly auditable closed-form decision function on three features (the ETH/BTC ratio, log returns at lag 6, and the 30-day natural gas return) with a 100\% symbolification rate (15 of 15 edges) and a post-symbolic test-fold accuracy of 51.85\%, a 5.18 percentage-point gain over the spline KAN's 46.67\%. The formula uses six smooth primitives ($\sin$, $\cos$, $\tanh$, $x^2$, $x^3$, $x^4$) and four outer terms with structurally distinct outer wrappers. All three surviving features have finite, well-defined gradients at the dataset median. Two structural findings from the cross-fold analysis support the symbolic-extraction reading. A lag feature (log returns at lag 6) survives extraction, validating the lag-features-in-MDA-pool design choice of Section 3.4.1. No on-chain feature survives extraction or reaches 50\% cross-fold selection frequency, so the Q1 ``on-chain free lunch'' hypothesis fails at the cross-fold stability level. The absolute post-symbolic accuracy of 51.85\% remains below the 56.85\% majority baseline, meaning the symbolic representation faithfully captures a model that did not extract directional signal on the deployment fold. Interpretability is therefore a separable contribution from predictive accuracy: closed-form transparency, smooth primitives, finite gradients, and three economically interpretable features are themselves the deliverable, and they do not wait on predictive significance to be valuable.

Beyond the current design, several methodological extensions are available. KAN 2.0 on the deployment fold \citep{liu_kan2_2024} would let multiplication nodes discover multiplicative feature interactions that the additive single-hidden-layer KAN of this thesis cannot represent, with the same four-step symbolic pipeline of Section 3.11 applied unchanged. Per-fold symbolic extraction across all 28 CPCV folds would let the analysis count which features and primitives recur and which are one-off, and is the highest-leverage methodological extension available. Higher-frequency (hourly) data would multiply CUSUM event counts by approximately 24 and unlock microstructure-feature spaces such as order-book imbalance and intraday volume profiles. Further extensions worth pursuing include regime-conditional analysis across BTC's distinct historical regimes, a meta-labelling layer \citep{singh_joubert_2019} above the bet-sizing stage, alternative assets such as ETH and gold, and a walk-forward versus CPCV comparison to isolate the contribution of the combinatorial purged design.

Three limitations of the current design themselves motivate future work. First, the results are conditional on a set of reasonable but untested labelling choices: the CUSUM threshold $h = 1.0\bar{\sigma}$, the $\pm 1.5\sigma_t$ profit and stop barriers with $\sigma_t$ from the span-50 EWMA, the 10-day vertical barrier, and the 2\% zero-class removal. A sensitivity analysis over this grid would quantify how much of the observed cross-model rank order is a property of the pipeline itself rather than of these specific constants. Second, the 73-feature universe carries natural redundancy across the technical, mathematical, and lag families; MDA selects the top-scoring features per fold but does not decorrelate them, so a hierarchical-clustering or variance-inflation pre-pass could sharpen the selected sets and clarify whether collinearity distorts the ranking. Third, the feature universe omits several families that recent crypto-quant work treats as first-class signals: investor attention (Google Trends, Wikipedia views), derivatives-market variables (perpetual funding rates, open interest, options-implied volatility), stablecoin supply dynamics, and BTC dominance. Adding these families under the same CPCV protocol would test whether the null result on trained models is a property of the feature universe or of BTC direction prediction more broadly.

Methodological honesty is not a constraint on financial machine learning. It is the prerequisite for treating interpretability as a contribution in its own right, and for reading buy-and-hold's dominance over every trained model as a result rather than a failure.